% Chapter 2: the literature the proposal builds on.
\chapter{Background and related work}\label{ch:background}

\section{Mapping urban temperature}

Urban climatology has long described the city as a mosaic of local
climate zones whose thermal behaviour follows from built form and land
cover \parencite{stewart2012local}. Statistical maps of air temperature
interpolate between stations using these covariates, most often through
land-use regression or kriging with external drift
\parencite{hjort2016spatial}. Recent work replaces linear trends with
gradient-boosted trees or neural networks, gaining accuracy but losing
the calibrated uncertainty that a placement rule needs
\parencite{venter2021hyperlocal}.

Gaussian processes keep that uncertainty explicit. With a kernel that
combines spatial distance and covariate similarity, they predict both a
mean temperature and a variance at every unobserved location
\parencite{rasmussen2006gaussian}. Sparse approximations make them
tractable for city-wide grids of a few hundred thousand cells
\parencite{hensman2013gaussian}.

\section{Optimal and adaptive sensor placement}

Choosing sites to maximise information is a classical problem in
spatial statistics. Mutual information between observed and unobserved
locations is submodular, so a greedy rule achieves at least
$(1-1/e)$ of the optimum \parencite{krause2008near}. Adaptive variants
place sensors one batch at a time and refit the model between batches,
which matters when the kernel itself is uncertain
\parencite{garnett2023bayesian}.

These methods treat all locations as equally important. Cost-sensitive
and task-driven variants exist in robotics and hydrology, where the
objective is a downstream decision rather than a map
\parencite{marchant2014bayesian}, but none has been applied to the
distribution of heat exposure across a population.

\section{The gap}

Three things are missing from current practice. No placement rule
weights uncertainty by exposure; no study tests whether a placement
policy learned in one city transfers to another; and no adaptive scheme
has been validated in a live urban deployment, where sensors fail and
access is negotiated site by site. The research questions in the next
chapter address each gap in turn.
